% !TeX root = ActuarialFormSheet_MBFA-es.tex
% !TeX spellcheck = es_ES
%\setlength{\parskip}{0cm} % paragraph spacing
%		\section*{Concepts de base}

\begin{f}[Definiciones]
	
	\textbf{Serie temporal} - sucesión de observaciones ordenadas en el tiempo con una frecuencia fija.
	
	Existen varias extensiones de las series temporales:
	\begin{itemize}[leftmargin=*]
		\item \textbf{Datos de panel} - consisten en una serie temporal para cada unidad de una sección transversal.
		\item \textbf{Secciones transversales agrupadas} - combinan secciones transversales correspondientes a distintos períodos.
	\end{itemize}
	
	\textbf{Proceso estocástico} - es una secuencia de variables aleatorias que están indexadas en el tiempo.
	
\end{f}
	
\begin{f}[Componentes de una serie temporal]{\ }
	
	\begin{itemize}[leftmargin=*]
		\item \textbf{Tendencia} - movimiento general a largo plazo de una serie.
		\item \textbf{Variaciones estacionales} - oscilaciones periódicas que se producen en un período igual o inferior a un año y que pueden identificarse fácilmente de un año a otro (generalmente debidas a factores climatológicos).
		\item \textbf{Variaciones cíclicas} - oscilaciones periódicas que se producen en un período superior a un año (son el resultado del ciclo económico).
		\item \textbf{Variaciones residuales} - movimientos que no siguen una oscilación periódica identificable (son el resultado de fenómenos puntuales).
	\end{itemize}
	
\end{f} 
\begin{f}[Tipos de modelos de series temporales]{\ }
	
	\begin{itemize}[leftmargin=*]
		\item \textbf{Modelos estáticos} - la relación entre \(y\) y \(x\) es contemporánea. Conceptualmente:
		\[
		y_{t} = \beta_{0} + \beta_{1} x_{t} + u_{t}
		\]
		\item \textbf{Modelos de retardos distribuidos} - la relación entre \(y\) y \(x\) no es contemporánea. Conceptualmente:
		\[
		y_{t} = \beta_{0} + \beta_{1} x_{t} + \beta_{2} x_{t-1} + \cdots + \beta_{s} x_{t-(s-1)} + u_{t}
		\]
		El efecto acumulado a largo plazo sobre \(y\) cuando \(\Delta x\) es:
		\[
		\beta_{1} + \beta_{2} + \cdots + \beta_{s}
		\]
		\item \textbf{Modelos dinámicos} - incluyen retardos de la variable dependiente (endogeneidad). Conceptualmente:
		\[
		y_{t} = \beta_{0} + \beta_{1} y_{t-1} + \cdots + \beta_{s} y_{t-s} + u_{t}
		\]
		\item \textbf{Combinaciones} de los modelos anteriores, como los modelos racionales de retardos distribuidos (retardos distribuidos + dinámica.
	\end{itemize}
	
\end{f}
 \hrule 
 
 \begin{f}[Hipótesis del modelo MCO en series temporales]
 	
 	Bajo estas hipótesis, el estimador MCO presenta buenas propiedades.
 	\textbf{Hipótesis de Gauss--Markov} extendidas a las series temporales:
 	
 	\begin{enumerate}[leftmargin=*, label=t\arabic{*}.]
 		
 		\item \textbf{Linealidad de los parámetros y dependencia débil}.
 		
 		\begin{enumerate}[leftmargin=*, label=\alph{*}.]
 			\item \(y_{t}\) debe ser una función lineal de los \(\beta\).
 			\item El proceso estocástico \(\lbrace(x_{t},y_{t}):t=1,2,\ldots,T\rbrace\) es estacionario y presenta dependencia débil.
 		\end{enumerate}
 		
 		\item \textbf{Ausencia de colinealidad perfecta}.
 		
 		\begin{itemize}[leftmargin=*]
 			\item Ninguna variable explicativa es constante: \(\Var(x_{j})\neq0,\;\forall j=1,\ldots,k\).
 			\item No existe una relación lineal exacta entre las variables explicativas.
 		\end{itemize}
 		
 		\item \textbf{Media condicional nula y ausencia de correlación}.
 		
 		\begin{enumerate}[leftmargin=*, label=\alph{*}.]
 		\item No existen errores sistemáticos: \(\E(u\mid x_{1},\ldots,x_{k})=\E(u)=0\) \(\rightarrow\) \textbf{exogeneidad fuerte} (a implica b).
 		\item No se omiten variables relevantes: \(\Cov(x_{j},u)=0,\;\forall j=1,\ldots,k\) \(\rightarrow\) \textbf{exogeneidad débil}.
 		\end{enumerate}
 		
 		\item \textbf{Homoscedasticidad}. La varianza de los residuos es constante para todo \(x\): \( \Var(u\mid x_{1},\ldots,x_{k})=\sigma^{2}_{u}\)
 		
 		\item \textbf{Ausencia de autocorrelación}. Los residuos no contienen información sobre los demás residuos: \( \Corr(u_t,u_s\mid x_{1},\ldots,x_{k})=0, \qquad \forall\, t\neq s\)
 		
 		\item \textbf{Normalidad}. Los residuos son independientes e idénticamente distribuidos (\textbf{i.i.d.}): \(u\sim\mathcal N(0,\sigma_u^2).\)
 		\item \textbf{Tamaño de la muestra}. El número de observaciones disponibles debe ser mayor que el número de parámetros a estimar \((k+1)\). (Esta condición se satisface automáticamente en un contexto asintótico.)
 		
 	\end{enumerate}
 	
 \end{f}
 
\begin{f}[Propiedades asintóticas del estimador MCO]
	
	Bajo las hipótesis del modelo econométrico y el teorema del límite central:
	
	\begin{itemize}[leftmargin=*]
		
		\item Si se verifican t1 a t3a: el estimador MCO es \textbf{insesgado}\(\E(\hat{\beta}_{j})=\beta_{j}.\)
		\item Si se verifican t1 a t3: el estimador MCO es \textbf{consistente}. \(\mathrm{plim}(\hat{\beta}_{j})=\beta_{j}.\) (Con t3b, t3a deja de cumplirse: existe exogeneidad débil, por lo que el estimador es sesgado pero consistente.) 
		\item Si se mantienen las hipótesis t1 a t5: el estimador MCO presenta \textbf{normalidad asintótica} (por tanto, t6 se satisface automáticamente cuando el tamaño muestral es suficientemente grande):
		\( \sqrt{T}\left(\hat{\beta}-\beta\right) \xrightarrow{d} \mathcal N\!\left(0,\Var(\hat{\beta})\right)\)
		\item Si se verifican t1 a t5: el estimador MCO es \textbf{eficiente} (teorema de Gauss--Markov), es decir, posee la menor varianza entre los estimadores lineales insesgados (\textbf{BLUE}).
	\end{itemize}
	
\end{f}  \hrule

\begin{f}[Tendencias y estacionalidad]
	
	\textbf{Regresión espuria} - se produce cuando la relación entre \(y\) y \(x\) se debe a factores que afectan a \(y\) y que están correlacionados con \(x\), \(\Corr(x_{j},u)\neq0\). Se trata del \textbf{incumplimiento de la hipótesis t3}.
	
	Dos series temporales pueden presentar la misma tendencia (o tendencias opuestas), lo que puede dar lugar a un elevado nivel de correlación. Esto puede generar una falsa apariencia de causalidad, conocida como \textbf{regresión espuria}. Dado el modelo:
	
	\begin{center}
		\(y_{t}=\beta_{0}+\beta_{1}x_{t}+u_{t}\)
	\end{center}
	
	donde:
	
	\begin{center}
		\(y_{t}=\alpha_{0}+\alpha_{1}\text{Tendencia}+v_{t}\)
		
		\(x_{t}=\gamma_{0}+\gamma_{1}\text{Tendencia}+v_{t}\)
	\end{center}
	
	La incorporación de una tendencia al modelo puede resolver este problema:
	
	\begin{center}
		\(y_{t}=\beta_{0}+\beta_{1}x_{t}+\beta_{2}\text{Tendencia}+u_{t}\)
	\end{center}
	
	La tendencia puede ser lineal o no lineal (cuadrática, cúbica, exponencial, etc.).
	
	Otro método consiste en utilizar el \textbf{filtro de Hodrick--Prescott} para extraer la tendencia y el componente cíclico.
	
\end{f}

\begin{f}[Estacionalidad]
	
	Una serie temporal puede presentar estacionalidad. Esto significa que la serie está sujeta a variaciones o patrones estacionales, normalmente relacionados con las condiciones climáticas.
	
	Por ejemplo, el PIB (en negro) suele ser mayor en verano y menor en invierno. La serie desestacionalizada (en naranja {\color{OrangeProfondIRA}}) se muestra para facilitar la comparación.
	
	\begin{tikzpicture}[xscale=0.4, yscale=.15]
		% \draw [step=1, gray, very thin] (0, 0) grid (20, 20);
		\draw [thick, <->] (0, 20) node [anchor=south] {\(y\)} -- (0, 0) -- (20, 0) node [anchor=south] {\(t\)};
		\draw [thick, black] 
		(0.0, 2.794) -- (0.5, 4.810) -- 
		(1.0, 2.500) -- (1.5, 7.619) -- 
		(2.0, 6.031) -- (2.5, 8.840) -- 
		(3.0, 5.420) -- (3.5, 10.855) -- 
		(4.0, 8.474) -- (4.5, 9.695) -- 
		(5.0, 5.481) -- (5.5, 9.512) -- 
		(6.0, 7.680) -- (6.5, 9.573) -- 
		(7.0, 5.787) -- (7.5, 10.366) -- 
		(8.0, 8.291) -- (8.5, 9.451) -- 
		(9.0, 5.604) -- (9.5, 10.099) -- 
		(10.0, 8.962) -- (10.5, 11.282) -- 
		(11.0, 7.130) -- (11.5, 11.709) -- 
		(12.0, 8.962) -- (12.5, 11.526) -- 
		(13.0, 8.168) -- (13.5, 13.358) -- 
		(14.0, 11.099) -- (14.5, 14.213) -- 
		(15.0, 10.916) -- (15.5, 16.290) -- 
		(16.0, 14.396) -- (16.5, 16.595) -- 
		(17.0, 13.419) -- (17.5, 18.000) -- 
		(18.0, 16.106) -- (18.5, 16.900); 
		\draw [thick, OrangeProfondIRA] 
		(0.0, 3.7939) -- (0.5, 3.9982) -- 
		(1.0, 3.9000) -- (1.5, 4.9183) -- 
		(2.0, 6.0905) -- (2.5, 6.9397) -- 
		(3.0, 6.9998) -- (3.5, 7.5450) -- 
		(4.0, 7.4733) -- (4.5, 7.6947) -- 
		(5.0, 7.4809) -- (5.5, 7.5115) -- 
		(6.0, 7.6794) -- (6.5, 7.5725) -- 
		(7.0, 7.7863) -- (7.5, 8.3336) -- 
		(8.0, 7.9901) -- (8.5, 8.1504) -- 
		(9.0, 8.6031) -- (9.5, 8.9008) -- 
		(10.0, 8.9618) -- (10.5, 8.7176) -- 
		(11.0, 8.9998) -- (11.5, 9.1901) -- 
		(12.0, 9.3618) -- (12.5, 9.3733) -- 
		(13.0, 10.6321) -- (13.5, 11.0588) --
		(14.0, 11.3992) -- (14.5, 12.2137) -- 
		(15.0, 12.5160) -- (15.5, 13.5901) -- 
		(16.0, 14.0969) -- (16.5, 15.2954) -- 
		(17.0, 15.3198) -- (17.5, 16.2000) -- 
		(18.0, 16.9069) -- (18.5, 17.2008);
	\end{tikzpicture}
	
	\begin{itemize}[leftmargin=*]
		\item Este problema da lugar a una \textbf{regresión espuria}. Un ajuste estacional puede resolverlo.
	\end{itemize}
	
	Un \textbf{ajuste estacional} sencillo consiste en crear variables ficticias estacionales e incorporarlas al modelo. Por ejemplo, para series trimestrales (\(Qq_t\) son variables binarias):
	
	\begin{center}
		\(y_{t}=\beta_{0}+\beta_{1}Q2_{t}+\beta_{2}Q3_{t}+\beta_{3}Q4_{t}+\beta_{4}x_{1t}+\cdots+\beta_{k}x_{kt}+u_{t}\)
	\end{center}
	
	Otro método consiste en desestacionalizar las variables y, a continuación, realizar la regresión con las variables ajustadas:
	
	\begin{center}
		\(z_{t}=\beta_{0}+\beta_{1}Q2_{t}+\beta_{2}Q3_{t}+\beta_{3}Q4_{t}+v_{t}\rightarrow\hat{v}_{t}+\E(z_{t})=\hat{z}_{t}^{sa}\)
		
		\(\hat{y}_{t}^{sa}=\beta_{0}+\beta_{1}\hat{x}_{1t}^{sa}+\cdots+\beta_{k}\hat{x}_{kt}^{sa}+u_{t}\)
	\end{center}
	
	Existen métodos mucho más eficaces y complejos para desestacionalizar una serie temporal, como el método \textbf{X-13ARIMA-SEATS}.
	
\end{f}

\begin{f}[Autocorrelación]
	
	El residuo de cualquier observación, \(u_t\), está correlacionado con el residuo de cualquier otra observación. Las observaciones no son independientes. Se trata de un caso de \textbf{incumplimiento} de la \textbf{hipótesis t5}.
	
	\begin{center}
		\(\Corr(u_{t},u_{s}\mid x_{1},\ldots,x_{k})=\Corr(u_{t},u_{s})\neq0,\;\forall\,t\neq s\)
	\end{center}
	
\end{f}    
\begin{f}[Consecuencias de la autocorrelación]
	
	La autocorrelación no introduce sesgo ni inconsistencia en los estimadores MCO si se cumple la hipótesis de exogeneidad estricta.
	
	Sin embargo:
	
	\begin{itemize}[leftmargin=*]
		\item los estimadores MCO dejan de ser eficientes (ya no son \textbf{BLUE});
		\item las varianzas estimadas son incorrectas;
		\item los errores estándar están sesgados;
		\item los contrastes \(t\) y \(F\), así como los intervalos de confianza, dejan de ser válidos.
	\end{itemize}
	
	En consecuencia, la inferencia estadística puede conducir a conclusiones erróneas.
	
\end{f}  

\begin{f}[Detección]
{\ }

\begin{itemize}[leftmargin=*]
	\item \textbf{Diagramas de dispersión}: busca patrones de dispersión de \(u_{t - 1}\) con respecto a \(u_{t}\).
	
	\setlength{\multicolsep}{0pt}
	\setlength{\columnsep}{6pt}
	\begin{multicols}{3}
		\begin{center}
			\begin{tikzpicture}[scale=0.11]
				\node at (16, 20) {\textbf{Ac.}}; 
				\draw [thick, ->] (0, 10) -- (20, 10) node [anchor=south] {\(u_{t - 1}\)}; 
				\draw [thick, -] (0, 0) -- (0, 20) node [anchor=west] {\(u_{t}\)}; 
				\draw plot [only marks, mark=*, mark size=6, domain=2:18, samples=50] (\x, {-0.2*(\x - 10)^2 + 13 + 6*rnd}); 
				\draw [thick, dashed, OrangeProfondIRA, -latex] plot [domain=2:18] (\x, {-0.2*(\x - 10)^2 + 16});
			\end{tikzpicture}
		\end{center}
		
		\columnbreak
		
		\begin{center}
			\begin{tikzpicture}[scale=0.11]
				\node at (16, 20) {\textbf{Ac. \(+\)}}; 
				\draw [thick, ->] (0, 10) -- (20, 10) node [anchor=north] {\(u_{t - 1}\)}; 
				\draw [thick, -] (0, 0) -- (0, 20) node [anchor=west] {\(u_{t}\)}; 
				\draw plot [only marks, mark=*, mark size=6, domain=2:18, samples=20] (\x, {5*rnd + 2.5 + 0.5*\x}); 
				\draw [thick, dashed, OrangeProfondIRA, -latex] plot [domain=2:18] (\x, {5 + 0.5*\x});
			\end{tikzpicture}
		\end{center}
		
		\columnbreak
		
		\begin{center}
			\begin{tikzpicture}[scale=0.11]
				\node at (16, 20) {\textbf{Ac. \(-\)}}; 
				\draw [thick, ->] (0, 10) -- (20, 10) node [anchor=south] {\(u_{t - 1}\)}; 
				\draw [thick, -] (0, 0) -- (0, 20) node [anchor=west] {\(u_{t}\)}; 
				\draw plot [only marks, mark=*, mark size=6, domain=2:18, samples=20] (\x, {5*rnd + 12.5 - 0.5*\x}); 
				\draw [thick, dashed, OrangeProfondIRA, -latex] plot [domain=2:18] (\x, {15 - 0.5*\x});
			\end{tikzpicture}
		\end{center}
	\end{multicols}
	
	\begin{multicols}{2}
		\item \textbf{Correlograma}: función de autocorrelación (ACF) y función de autocorrelación parcial (PACF).
	
	\columnbreak
	
	\begin{itemize}[leftmargin=*]
		\item Eje Y: correlación.
		\item Eje X: número de desplazamientos.
		\item Zona gris: \(\pm 1,96/T^{0,5}\)
	\end{itemize}
\end{multicols}

\begin{center}
\begin{tikzpicture}[scale=0.25]
% acf plot
\node at (-2.5, 14) {\small \rotatebox{90}{\textbf{ACF}}}; 
\node at (-1, 17.5) {\small 1};
\node at (-1, 14) {\small 0};
\node at (-1, 10.5) {\small -1}; 
\fill [lightgray] (0, 13) rectangle (30.5, 15); 
\draw [dashed, thin] (0, 14) -- (30.5, 14); 
\draw [thick, |->] (0, 18) -- (0, 10) -- (30.5, 10);
\fill [OrangeProfondIRA] (2, 14) rectangle (2.5, 17.95);
\fill [OrangeProfondIRA] (5, 14) rectangle (5.5, 16.96);
\fill [OrangeProfondIRA] (8, 14) rectangle (8.5, 16.22); 
\fill [OrangeProfondIRA] (11, 14) rectangle (11.5, 15.67);
\fill [OrangeProfondIRA] (14, 14) rectangle (14.5, 15.25);
\fill [OrangeProfondIRA] (17, 14) rectangle (17.5, 14.94);
\fill [OrangeProfondIRA] (20, 14) rectangle (20.5, 14.70);
\fill [OrangeProfondIRA] (23, 14) rectangle (23.5, 14.53);
\fill [OrangeProfondIRA] (26, 14) rectangle (26.5, 14.40);
\fill [OrangeProfondIRA] (29, 14) rectangle (29.5, 14.30);
% pacf plot
\node at (-2.5, 4) {\small \rotatebox{90}{\textbf{PACF}}};
\node at (-1, 7.5) {\small 1};
\node at (-1, 4) {\small 0};
\node at (-1, 0.5) {\small -1};
\fill [lightgray] (0, 3) rectangle (30.5, 5);
\draw [dashed, thin] (0, 4) -- (30.5, 4);
\draw [thick, |->] (0, 8) -- (0, 0) -- (30.5, 0);
\fill [OrangeProfondIRA] (2, 4) rectangle (2.5, 7.90);
\fill [OrangeProfondIRA] (5, 4) rectangle (5.5, 7.00);
\fill [OrangeProfondIRA] (8, 4) rectangle (8.5, 3.47);
\fill [OrangeProfondIRA]	(11, 4) rectangle (11.5, 4.24);
\fill [OrangeProfondIRA] (14, 4) rectangle (14.5, 4.43);
\fill [OrangeProfondIRA] (17, 4) rectangle (17.5, 4.89);
\fill [OrangeProfondIRA] (20, 4) rectangle (20.5, 3.09);
\fill [OrangeProfondIRA] (23, 4) rectangle (23.5, 3.58);
\fill [OrangeProfondIRA] (26, 4) rectangle (26.5, 4.46);
\fill [OrangeProfondIRA] (29, 4) rectangle (29.5, 4.86);
\end{tikzpicture}
\end{center}
	

\textbf{Proceso MA(\(q\))}. \underline{FAC}: únicamente los primeros \(q\) coeficientes son significativos; los demás se anulan bruscamente. \underline{FACP}: presenta una rápida disminución exponencial amortiguada o un patrón de ondas sinusoidales.

\textbf{Proceso AR(\(p\))}. \underline{FAC}: presenta una rápida disminución exponencial amortiguada o un patrón de ondas sinusoidales. \underline{FACP}: únicamente los primeros \(p\) coeficientes son significativos; los demás se anulan bruscamente.

\textbf{Proceso ARMA(\(p,q\))}. \underline{FAC} y \underline{FACP}: los coeficientes no se anulan bruscamente y presentan una rápida disminución.

Si los coeficientes de la FAC no disminuyen rápidamente, ello indica claramente una falta de estacionariedad en la media.

\item \textbf{Contrastes formales} - En general, \(H_{0}\): ausencia de autocorrelación.

Suponiendo que \(u_{t}\) sigue un proceso AR(1):

\begin{center}
	\(u_{t}=\rho_{1}u_{t-1}+\varepsilon_{t}\)
\end{center}

donde \(\varepsilon_{t}\) es un ruido blanco.

\textbf{Contraste \(t\) para AR(1)} (regresores exógenos):

\begin{center}
	\(t=\dfrac{\hat{\rho}_{1}}{\se(\hat{\rho}_{1})}\sim t_{T-k-1,\alpha/2}\)
\end{center}

\begin{itemize}[leftmargin=*]
	\item \(H_{1}\): autocorrelación de primer orden, AR(1).
\end{itemize}

\textbf{Estadístico de Durbin--Watson} (regresores exógenos y normalidad de los residuos):

\begin{center}
	\(d=\dfrac{\sum_{t=2}^{n}(\hat{u}_{t}-\hat{u}_{t-1})^{2}}{\sum_{t=1}^{n}\hat{u}_{t}^{2}}\approx2\cdot(1-\hat{\rho}_{1})\)
\end{center}

Donde \(0\leq d\leq4\).

\begin{itemize}[leftmargin=*]
	\item \(H_{1}\): autocorrelación de primer orden, \(AR(1)\).
\end{itemize}
	
\begin{center}
			\begin{tabular}{ c | c | c | c }
			\(d =\)          & 0 & 2 & 4  \\ \hline
			\(\rho \approx\) & 1 & 0 & -1
		\end{tabular}

\end{center}		

	\begin{tikzpicture}[scale=0.3]
		\fill [lightgray] (5, 0) rectangle (9, 6); 
		\draw (5, 0) -- (5, 6);
		\draw (9, 0) -- (9, 6);
		\fill [lightgray] (16, 0) rectangle (20, 6);
		\draw (16, 0) -- (16, 6);
		\draw (20, 0) -- (20, 6);
		\draw [thick] (0, 6) -- (0, 0) -- (25, 0);
		\draw [dashed] (12.5, 0) -- (12.5, 6);
		\node at (-0.5, 6.5) {\small \(f(d)\)};
		\node at (0, -0.6) {\small 0};
		\node at (5, -0.6) {\small \(d_{L}\)};
		\node at (9, -0.6) {\small \(d_{U}\)};
		\node at (12.5, -0.6) {\small 2};
		\node at (16.7, -0.6) {\tiny \((4 - d_{U})\)};
		\node at (20.7, -0.6) {\tiny \((4 - d_{L})\)};
		\node at (25, -0.6) {\small 4};
		\node at (2.5, 3.5) {\small Rej. \(H_{0}\)};
		\node at (2.5, 2.5) {\small AR \(+\)};
		\node [text=OrangeProfondIRA] at (7, 3) {\textbf{?}};
		\node at (12.5, 3.5) {\small Not rej. \(H_{0}\)};
		\node at (12.5, 2.5) {\small No AR};
		\node [text=OrangeProfondIRA] at (18, 3) {\textbf{?}};
		\node at (22.5, 3.5) {\small Rej. \(H_{0}\)};
		\node at (22.5, 2.5) {\small AR \(-\)};
	\end{tikzpicture}

	\textbf{h de Durbin} (regresores endógenos):

\begin{center}
	\(h = \hat{\rho} \cdot \sqrt{\dfrac{T}{1 - T \cdot \upsilon}}\)
\end{center}

donde \(\upsilon\) es la varianza estimada del coeficiente asociado a la variable endógena.

\begin{itemize}[leftmargin=*]
	\item \(H_{1}\): autocorrelación de primer orden, AR(1).
\end{itemize}

\textbf{Contraste de Breusch--Godfrey} (regresores endógenos): permite detectar procesos MA(\(q\)) y AR(\(p\)) (\(\varepsilon_{t}\) es un ruido blanco):

\begin{itemize}[leftmargin=*]
	\item MA(\(q\)): \(u_{t}=\varepsilon_{t}-m_{1}u_{t-1}-\cdots-m_{q}u_{t-q}\)
	\item AR(\(p\)): \(u_{t}=\rho_{1}u_{t-1}+\cdots+\rho_{p}u_{t-p}+\varepsilon_{t}\)
\end{itemize}

Bajo \(H_{0}\): ausencia de autocorrelación:

\begin{center}
	\(\hfill T\cdot R^{2}_{\hat{u}_t}\underset{a}{\sim}\chi^{2}_{q} \hfill \textbf{o} \hfill T\cdot R^{2}_{\hat{u}_t}\underset{a}{\sim}\chi^{2}_{p} \hfill\)
\end{center}

\begin{itemize}[leftmargin=*]
	\item \(H_{1}\): autocorrelación de orden \(q\) (o \(p\)).
\end{itemize}

\textbf{Contraste \(Q\) de Ljung--Box}:

\begin{itemize}[leftmargin=*]
	\item \(H_{1}\): autocorrelación hasta el retardo \(h\).
\end{itemize}

\end{itemize}


\end{f}  

\begin{f}[Corrección]
	
	\begin{itemize}[leftmargin=*]
		\item Utilizar el método de \textbf{mínimos cuadrados ordinarios} (MCO) con un estimador de la matriz de varianzas-covarianzas \textbf{robusto frente a la heterocedasticidad y la autocorrelación} (HAC), como el propuesto por \textbf{Newey-West}.
		
		\item Utilizar los \textbf{mínimos cuadrados generalizados} (MCG). Supóngase que \(y_{t} = \beta_{0} + \beta_{1} x_{t} + u_{t}\), con \(u_{t} = \rho u_{t - 1}+ \varepsilon_{t}\), donde \(\lvert \rho \rvert < 1\) y \(\varepsilon_{t}\) es un \underline{ruido blanco}.
		
		\begin{itemize}[leftmargin=*]
			
			\item Si \(\rho\) es \textbf{conocido}, utilizar un \textbf{modelo cuasi-diferenciado}:
			
			\begin{center}
				\(y_{t} - \rho y_{t - 1}= \beta_{0} (1 - \rho) + \beta_{1} (x_{t} - \rho x_{t - 1}) + u_{t} - \rho u_{t - 1}\)
				
				\(y_{t}^{*} = \beta_{0}^{*} + \beta_{1}' x_{t}^{*} + \varepsilon_{t}\)
			\end{center}
			
			donde \(\beta_{1}' = \beta_{1}\); y estimarlo mediante MCO.
			
			\item Si \(\rho\) \textbf{no es conocido}, estimarlo, por ejemplo, mediante el \textbf{método iterativo de Cochrane-Orcutt} (el método de Prais-Winsten también es válido):
			
			\begin{enumerate}[leftmargin=*]
				\item Obtener \(\hat{u}_{t}\) a partir del modelo original.
				\item Estimar \(\hat{u}_{t} = \rho \hat{u}_{t-1} + \varepsilon_{t}\) y obtener \(\hat{\rho}\).
				\item Construir un modelo cuasi-diferenciado:
				
				\begin{center}
					\(y_{t} - \hat{\rho}y_{t - 1} = \beta_{0} (1 - \hat{\rho}) + \beta_{1} (x_{t} - \hat{\rho} x_{t - 1}) + u_{t} - \hat{\rho}u_{t - 1}\)
					
					\(y_{t}^{*} = \beta_{0}^{*} + \beta_{1}' x_{t}^{*} + \varepsilon_{t}\)
				\end{center}
				
				donde \(\beta_{1}' = \beta_{1}\); y estimarlo mediante MCO.
				
				\item Obtener \(\hat{u}_{t}^{*} = y_{t} - (\hat{\beta}_{0}^{*} + \hat{\beta}_{1}' x_{t}) \neq y_{t} - (\hat{\beta}_{0}^{*} + \hat{\beta}_{1}' x_{t}^{*})\).
				
				\item Repetir a partir del paso 2. El algoritmo finaliza cuando los parámetros estimados apenas varían entre iteraciones.
			\end{enumerate}
			
		\end{itemize}
		
		\item Si el problema persiste, buscar una \textbf{dependencia fuerte} en la serie.
		
	\end{itemize}
	
\end{f}

\columnbreak

\hrule

\begin{f}[Suavización exponencial]
	
	\begin{center}
		\(f_{t} = \alpha y_{t} + (1 - \alpha) f_{t - 1}\)
	\end{center}
	
	donde \(0 < \alpha < 1\) es el parámetro de suavización.
	
\end{f}  \hrule

\begin{f}[Predicción]
	
	Existen dos tipos de predicción:
	
	\begin{itemize}[leftmargin=*]
		\item Del valor medio de \(y\) para un valor específico de \(x\).
		\item De un valor individual de \(y\) para un valor específico de \(x\).
	\end{itemize}
	
	\textbf{Estadístico U de Theil} - compara los valores predichos con los obtenidos mediante una predicción ingenua basada en un mínimo de datos históricos.
	
	\begin{center}
		\(U = \sqrt{\frac{\sum_{t=1}^{T-1} \left( \frac{\hat{y}_{t+1} - y_{t+1}}{y_t} \right)^2}{\sum_{t=1}^{T-1} \left( \frac{y_{t+1} - y_t}{y_t} \right)^2}}\)
	\end{center}
	
	\begin{itemize}[leftmargin=*]
		\item \(<1\) : la predicción es mejor que la predicción ingenua.
		\item \(=1\) : la predicción es aproximadamente tan buena como la predicción ingenua.
		\item \(>1\) : la predicción es peor que la predicción ingenua.
	\end{itemize}
	
\end{f}  \hrule

\begin{f}[Estacionariedad]
	
	La estacionariedad permite identificar correctamente las relaciones entre las variables que permanecen constantes a lo largo del tiempo.
	
	\begin{itemize}[leftmargin=*]
		\item \textbf{Proceso estacionario} (estacionariedad estricta): si un conjunto de variables aleatorias se desplaza \(h\) períodos (cambios en el tiempo), la distribución conjunta de probabilidad debe permanecer inalterada.
		
		\item \textbf{Proceso no estacionario}: por ejemplo, una serie con tendencia, donde al menos la media cambia con el tiempo.
		
		\item \textbf{Proceso estacionario en covarianza} - constituye una forma más débil de estacionariedad:
		
		\begin{itemize}[leftmargin=*]
			\begin{multicols}{2}
				\item \(\E(x_{t})\) es constante.
				
				\columnbreak
				
				\item \(\Var(x_{t})\) es constante.
			\end{multicols}
			
			\item Para todo \(t\), \(h \geq 1\), \(\Cov(x_{t},x_{t+h})\) depende únicamente de \(h\), y no de \(t\).
		\end{itemize}
	\end{itemize}
	
\end{f}  \hrule

\begin{f}[Dependencia débil]
	
	La dependencia débil sustituye la hipótesis de muestreo aleatorio en las series temporales.
	
	\begin{itemize}[leftmargin=*]
		\item Un proceso estacionario \(\lbrace x_{t} \rbrace\) presenta \textbf{dependencia débil} cuando \(x_{t}\) y \(x_{t+h}\) son prácticamente independientes a medida que \(h\) tiende a infinito.
		
		\item Un proceso estacionario en covarianza presenta \textbf{dependencia débil} si la correlación entre \(x_{t}\) y \(x_{t+h}\) converge a \(0\) con suficiente rapidez cuando \(h \rightarrow \infty\) (son asintóticamente incorrelados).
	\end{itemize}
	
	Los procesos con dependencia débil se denominan \textbf{integrados de orden cero}, I(0). Algunos ejemplos son:
	
	\begin{itemize}[leftmargin=*]
		
		\item \textbf{Media móvil} - \(\lbrace x_{t} \rbrace\) es una media móvil de orden \(q\), MA(\(q\)):
		
		\begin{center}
			\(x_{t}=e_{t}+m_{1}e_{t-1}+\cdots+m_{q}e_{t-q}\)
		\end{center}
		
		donde \(\lbrace e_{t}:t=0,1,\ldots,T\rbrace\) es una secuencia \textsl{i.i.d.} con media nula y varianza \(\sigma_{e}^{2}\).
		
		\item \textbf{Proceso autorregresivo} - \(\lbrace x_{t} \rbrace\) es un proceso autorregresivo de orden \(p\), AR(\(p\)):
		
		\begin{center}
			\(x_{t}=\rho_{1}x_{t-1}+\cdots+\rho_{p}x_{t-p}+e_{t}\)
		\end{center}
		
		donde \(\lbrace e_{t}:t=1,2,\ldots,T\rbrace\) es una secuencia \textsl{i.i.d.} con media nula y varianza \(\sigma_{e}^{2}\).
		
		\textbf{Condición de estabilidad}: si \(1-\rho_{1}z-\cdots-\rho_{p}z^{p}=0\) para \(\lvert z\rvert>1\), entonces \(\lbrace x_{t}\rbrace\) es un proceso AR(\(p\)) estable con dependencia débil. Para un AR(1), la condición es: \(\lvert\rho_{1}\rvert<1\).
		
		\item \textbf{Proceso ARMA} - combinación de un proceso AR(\(p\)) y un proceso MA(\(q\)); \(\lbrace x_{t}\rbrace\) es un proceso ARMA(\(p,q\)):
		
		\begin{center}
			\(x_{t}=e_{t}+m_{1}e_{t-1}+\cdots+m_{q}e_{t-q}+\rho_{1}x_{t-1}+\cdots+\rho_{p}x_{t-p}\)
		\end{center}
		
	\end{itemize}
	
\end{f}  \hrule

 \begin{f}[Raíces unitarias]
 	
 	Un proceso es I(\(d\)), es decir, integrado de orden \(d\), si al diferenciarlo \(d\) veces se obtiene un proceso estacionario.
 	
 	Cuando \(d \geq 1\), el proceso se denomina \textbf{proceso con raíz unitaria}, o se dice que presenta una raíz unitaria.
 	
 	Un proceso presenta una raíz unitaria cuando no se cumple la condición de estabilidad (existen raíces sobre el círculo unitario).
 	
 \end{f}  \hrule
 
\begin{f}[Fuerte dependencia]
	
	En la mayoría de los casos, las series económicas presentan una fuerte dependencia (o son muy persistentes). Algunos ejemplos de \textbf{raíz unitaria} I(1):
	
	\begin{itemize}[leftmargin=*]
		\item \textbf{Marcha aleatoria}: un proceso AR(1) con \(\rho_{1} = 1\).
		
		\begin{center}
			\(y_{t} = y_{t - 1} + e_{t}\)
		\end{center}
		
		donde \(\lbrace e_{t} : t = 1, 2, \ldots, T \rbrace\) es una secuencia \textsl{i.i.d.} con media nula y varianza \(\sigma^{2}_{e}\).
		
		\item \textbf{Marcha aleatoria con deriva}: un proceso AR(1) con \(\rho_{1} = 1\) y una constante.
		
		\begin{center}
			\(y_{t} = \beta_{0} + y_{t - 1} + e_{t}\)
		\end{center}
		
		donde \(\lbrace e_{t} : t = 1, 2, \ldots, T \rbrace\) es una secuencia \textsl{i.i.d.} con media nula y varianza \(\sigma^{2}_{e}\).
	\end{itemize}
	
\end{f} 

\begin{f}[Pruebas de raíz unitaria]{\ }
	
	\begin{center}
		\begin{tabular}{ c | c | c }
			Prueba            & \(H_{0}\)    & Rechazar \(H_{0}\)                     \\ \hline
			ADF             & I(1)       & tau \textless \, Valor crítico    \\ \hline
			KPSS            & Nivel I(0) & mu \textgreater \, Valor crítico  \\
			& Tendencia I(0) & tau \textgreater \, Valor crítico \\ \hline
			Phillips-Perron & I(1)       & Z-tau \textless \, valor crítico  \\ \hline
			Zivot-Andrews   & I(1)       & tau \textless \, valor crítico
		\end{tabular}
	\end{center}
	\medskip
	
\end{f}  

\begin{f}[De la raíz unitaria a la dependencia débil]
	
	Un proceso integrado de orden \textbf{uno}, I(1), significa que \textbf{la primera diferencia} del proceso es \textbf{débilmente dependiente} o I(0) (y, por lo general, estacionaria). Por ejemplo, sea \(\lbrace y_{t} \rbrace\) un proceso aleatorio:
	
	\begin{multicols}{2}
		\begin{center}
			\(\Delta y_{t} = y_{t} - y_{t - 1} = e_{t}\)
		\end{center}
		
		donde \(\lbrace e_{t} \rbrace = \lbrace \Delta y_{t} \rbrace\) es \textsl{i.i.d.} \\
		
		Nota:
		\begin{itemize}[leftmargin=*]
			\item La primera diferencia de una serie elimina su tendencia.
			\item Los logaritmos de una serie estabilizan su varianza.
		\end{itemize}
		
		\columnbreak
		
			\begin{tikzpicture}[scale=0.18]
	% \draw [step=1, gray, very thin] (0, 0) grid (20, 20); 
	\draw [thick, <->] (0, 20) node [anchor=south west] {\(y, {\color{OrangeProfondIRA} \Delta y}\)} -- (0, 0) -- (20, 0) node [anchor=south] {\(t\)}; 
	\draw [thick, black] 
	(0.0, 2.000) -- (0.5, 2.459) -- 
	(1.0, 2.716) -- (1.5, 3.205) -- 
	(2.0, 3.571) -- (2.5, 3.952) -- 
	(3.0, 4.047) -- (3.5, 4.514) -- 
	(4.0, 4.719) -- (4.5, 5.160) -- 
	(5.0, 5.674) -- (5.5, 5.987) -- 
	(6.0, 6.242) -- (6.5, 6.471) -- 
	(7.0, 6.944) -- (7.5, 7.104) -- 
	(8.0, 7.584) -- (8.5, 8.087) -- 
	(9.0, 8.112) -- (9.5, 8.834) -- 
	(10.0, 9.470) -- (10.5, 9.718) -- 
	(11.0, 10.032) -- (11.5, 10.491) -- 
	(12.0, 10.748) -- (12.5, 10.805) -- 
	(13.0, 11.016) -- (13.5, 11.439) --
	(14.0, 11.810) -- (14.5, 12.247) -- 
	(15.0, 12.668) -- (15.5, 13.052) -- 
	(16.0, 13.586) -- (16.5, 14.322) -- 
	(17.0, 14.913) -- (17.5, 15.704) -- 
	(18.0, 16.081) -- (18.5, 16.431); 
	\draw [thick, OrangeProfondIRA] 
	(0.5, 11.283) -- (1.0, 7.201) -- 
	(1.5, 11.889) -- (2.0, 9.405) -- 
	(2.5, 9.701) -- (3.0, 3.926) -- 
	(3.5, 11.454) -- (4.0, 6.136) -- 
	(4.5, 10.926) -- (5.0, 12.393) -- 
	(5.5, 8.345) -- (6.0, 7.157) -- 
	(6.5, 6.627) -- (7.0, 11.572) -- 
	(7.5, 5.235) -- (8.0, 11.703) -- 
	(8.5, 12.186) -- (9.0, 2.513) -- 
	(9.5, 16.607) -- (10.0, 14.869) -- 
	(10.5, 7.015) -- (11.0, 8.368) -- 
	(11.5, 11.283) -- (12.0, 7.196) -- 
	(12.5, 3.153) -- (13.0, 6.277) -- 
	(13.5, 10.547) -- (14.0, 9.517) -- 
	(14.5, 10.834) -- (15.0, 10.526) -- 
	(15.5, 9.754) -- (16.0, 12.816) -- 
	(16.5, 16.875) -- (17.0, 13.961) -- 
	(17.5, 18.000) -- (18.0, 9.644) -- 
	(18.5, 9.075);
\end{tikzpicture}
\end{multicols}
	
	
	
	\subsubsection*{De la raíz unitaria al porcentaje de variación}
	
	Cuando una serie I(1) es estrictamente positiva, se suele convertir a logaritmos antes de calcular la primera diferencia para obtener el porcentaje de variación (aproximado) de la serie:
	
	\begin{center}
		\(\Delta \log(y_{t}) = \log(y_{t}) - \log(y_{t - 1}) \approx \dfrac{y_t - y_{t - 1}} {y_{t - 1}}\)
	\end{center}
	
\end{f}  \hrule

\begin{f}[Cointegración]
	
	Cuando \textbf{dos series son I(1), pero una combinación lineal de ambas es I(0)}. En este caso, la regresión de una serie sobre la otra no es engañosa, sino que expresa algo sobre la relación a largo plazo. Se dice que las variables están cointegradas si tienen una tendencia estocástica común.
	
	Por ejemplo, \(\lbrace x_{t} \rbrace\) y \(\lbrace y_{t} \rbrace\) son I(1), pero \(y_{t} - \beta x_{t} = u_{t}\), donde \(\lbrace u_{t} \rbrace\) es I(0). (\(\beta\) es el parámetro de cointegración).
	
\end{f}  

\begin{f}[Prueba de cointegración]
	
	Siguiendo el ejemplo anterior:
	
	\begin{enumerate}[leftmargin=*]
		\item Estimar \(y_{t} = \alpha + \beta x_{t} + \varepsilon_{t}\) y obtener \(\hat{\varepsilon}_{t}\).
		\item Realizar una prueba ADF sobre \(\hat{\varepsilon}_{t}\) con una distribución modificada.
		
		El resultado de esta prueba es equivalente a:
		
		\begin{itemize}[leftmargin=*]
			\item \(H_{0}\): \(\beta = 0\) (sin cointegración)
			\item \(H_{1}\): \(\beta \neq 0\) (cointegración)
		\end{itemize}
		
		si la estadística de la prueba es \(>\) el valor crítico, rechazar \(H_0\).
	\end{enumerate}
	
\end{f}  \hrule

\begin{f}[Heterocedasticidad en las series temporales]
	
	La \textbf{hipótesis} afectada es \textbf{t4}, lo que conduce a \textbf{una ineficacia de los MCO}.
	
	Utilice pruebas como las de Breusch-Pagan o White, donde \(H_{0}\): no hay heterocedasticidad. Es \textbf{importante} para que las pruebas funcionen que no haya \textbf{autocorrelación}.
	
\end{f}  \hrule  
\begin{f}[ARCH]
	
	La heterocedasticidad autorregresiva condicional (ARCH) es un modelo que permite analizar una forma de heterocedasticidad dinámica, en la que la varianza del error sigue un proceso AR(\(p\)).
	
	Dado el modelo: \(y_{t} = \beta_{0} + \beta_{1} z_{t} + u_{t}\), donde hay AR(1) y heterocedasticidad:
	
	\begin{center}
		\(\E(u^{2}_{t} \mid u_{t - 1}) = \alpha_{0} + \alpha_{1} u^{2}_{t - 1}\)
	\end{center}
	
\end{f}  \hrule 
\begin{f}[GARCH]
	
	Un modelo general de heterocedasticidad autorregresiva condicional (GARCH) es similar al modelo ARCH, pero en este caso la varianza del error sigue un proceso ARMA(\(p, q\)).
\end{f}